\documentclass{article}
\usepackage{amsmath,amssymb}
\begin{document}

\section*{Qualitative change of motion: bound $\leftrightarrow$ unbound}

\subsection*{1. Classification and the conservation obstruction}
Relative motion of two bodies (reduced mass $\mu$, $V=-k/r$, with
$k=Gm_1m_2$ or $k=e^2/4\pi\epsilon_0$):
\begin{equation}
E=\tfrac12\mu\dot{\mathbf r}^{2}-\frac{k}{r},\qquad
L=\mu\,|\mathbf r\times\dot{\mathbf r}|,\qquad
\varepsilon=\sqrt{1+\frac{2EL^{2}}{\mu k^{2}}}
\end{equation}
\begin{itemize}
  \item $E<0$ ($\varepsilon<1$): ellipse, bound, compact energy surface;
  \item $E=0$ ($\varepsilon=1$): parabola, the separatrix;
  \item $E>0$ ($\varepsilon>1$): hyperbola, scattering, non-compact energy surface.
\end{itemize}
Since $\dot E=\dot L=0$, the sets $\{E<0\}$ and $\{E\ge0\}$ are invariant:
\emph{an isolated two-body system can neither capture nor ionize.}

\subsection*{2. Transition = coupling that changes $E$}
Add an external force $\mathbf F_{\rm ext}$ (or explicit time dependence):
\begin{equation}
\mu\ddot{\mathbf r}=-\nabla V+\mathbf F_{\rm ext},\qquad
\dot E=\frac{\partial V}{\partial t}+\mathbf F_{\rm ext}\!\cdot\dot{\mathbf r}.
\end{equation}
Typical sources:
\begin{align}
\text{third body (tidal):}\quad
&\mathbf F_{\rm ext}=\mu\,\frac{GM_3}{D^{3}}
\left[3\hat{\mathbf D}(\hat{\mathbf D}\!\cdot\mathbf r)-\mathbf r\right]\\
\text{drag / tidal friction:}\quad
&\mathbf F_{\rm ext}=-\gamma\,\dot{\mathbf r}\ \Rightarrow\ \dot E=-\gamma\dot{\mathbf r}^{2}<0\\
\text{radiation reaction (Larmor):}\quad
&\dot E=-\frac{2}{3}\,\frac{k\,a^{2}}{c^{3}},\qquad a=\frac{k}{m r^{2}}\\
\text{growing central mass:}\quad
&\dot E=\partial_tV=-\frac{G\dot M m}{r}<0
\end{align}
Capture of a body arriving with relative speed $v_\infty$ requires
\begin{equation}
\Delta E=\int \mathbf F_{\rm ext}\!\cdot\dot{\mathbf r}\,dt<-\tfrac12\mu v_\infty^{2},
\end{equation}
and ionization/escape requires $\Delta E\ge|E|$.

\subsection*{3. Planet captures satellite: restricted three-body (Hill) problem}
Rotating frame centred on the planet ($M_p$), mean motion $n$ of the planet about the Sun:
\begin{align}
\ddot x-2n\dot y&=3n^{2}x-\frac{GM_p\,x}{r^{3}},\\
\ddot y+2n\dot x&=-\frac{GM_p\,y}{r^{3}},\\
\ddot z&=-n^{2}z-\frac{GM_p\,z}{r^{3}}.
\end{align}
Conserved Jacobi constant and Hill radius:
\begin{equation}
C_J=3n^{2}x^{2}-n^{2}z^{2}+\frac{2GM_p}{r}-v^{2},\qquad
r_H=\left(\frac{GM_p}{3n^{2}}\right)^{1/3}.
\end{equation}
Since $v^{2}\ge0$, motion is confined to
$3n^{2}x^{2}-n^{2}z^{2}+2GM_p/r\ge C_J$. The necks at $L_{1,2}$ open when
$C_J<C_{L}=9n^{2}r_H^{2}$.
\begin{itemize}
  \item With $C_J$ conserved and the dynamics time-reversible, a body that enters
        the Hill region also leaves it: permanent capture of a scattering orbit has
        measure zero, and \emph{temporary capture} is a long sojourn on a chaotic set.
  \item Permanent capture needs a change $\Delta C_J$: gas drag, tidal
        dissipation, planet mass growth, collision, or an exchange reaction
        (binary--planet encounter, one partner ejected, carrying away the energy).
\end{itemize}

\subsection*{4. Proton + electron $\leftrightarrow$ hydrogen}
\begin{equation}
\mathrm p+\mathrm e^-(E_k>0)\;\longrightarrow\;\mathrm H(nlm)+\gamma,
\qquad
\hbar\omega=\tfrac12mv^{2}+\frac{I}{n^{2}},\quad I=\frac{mk^{2}}{2\hbar^{2}},
\end{equation}
\begin{equation}
\mathrm H(nlm)+\gamma\;\longrightarrow\;\mathrm p+\mathrm e^-
\quad\text{(photoionization, the reverse process).}
\end{equation}
The closed system (e + p + field) is Hamiltonian; the photon carries away the
energy, so the subsystem (e, p) changes from $E_{12}>0$ to $E_{12}<0$.

Quantum description: the spectrum is discrete for $E<0$ and continuous for $E>0$, and the transition is a
matrix element between a Coulomb continuum state $|\mathbf k\rangle$ and a bound state:
\begin{equation}
c_f(t)=-\frac{i}{\hbar}\int_0^{t}\langle f|H_{\rm int}(t')|i\rangle\,
e^{i\omega_{fi}t'}\,dt',\qquad
H_{\rm int}=\frac{e}{m}\,\mathbf A\!\cdot\!\mathbf p,
\end{equation}
\begin{equation}
\Gamma_{i\to f}=\frac{2\pi}{\hbar}\,|\langle f|H_{\rm int}|i\rangle|^{2}\,\rho(E_f),
\qquad
A_{i\to f}=\frac{4k\,\omega^{3}}{3\hbar c^{3}}\,|\langle f|\mathbf r|i\rangle|^{2}
\ \ \text{(dipole approx.)}
\end{equation}
Detailed balance (Milne relation) links the two directions:
\begin{equation}
g_\gamma g_n\Big(\frac{\hbar\omega}{c}\Big)^{2}\sigma_{\rm ion}
= g_e g_p\,(mv)^{2}\,\sigma_{\rm rec}
\;\Rightarrow\;
\sigma_{\rm rec}=\frac{g_\gamma g_n}{g_e g_p}\,
\frac{(\hbar\omega)^{2}}{m^{2}c^{2}v^{2}}\,\sigma_{\rm ion},
\end{equation}
with $g_\gamma=2$, $g_e=2$, $g_p=1$, $g_n=2n^2$.

Population dynamics and thermal equilibrium (Saha):
\begin{equation}
\frac{dn_{\rm H}}{dt}=\alpha(T)\,n_en_p-\Gamma_{\rm ion}n_{\rm H},\quad
\alpha=\langle\sigma_{\rm rec}v\rangle,\qquad
\frac{n_en_p}{n_{\rm H}}=\frac{2g_p}{g_n}
\Big(\frac{2\pi mk_BT}{h^{2}}\Big)^{3/2}e^{-I/k_BT}.
\end{equation}

\subsection*{5. Unified criterion}
Define $E_{12}(t)=H_{12}$ for the subsystem. The state is bound iff $E_{12}<0$;
capture and ionization/escape are the crossings $E_{12}=0$ with $\dot E_{12}<0$ and $\dot E_{12}>0$.
Classically, $\Delta E=\int\mathbf F_{\rm ext}\!\cdot\dot{\mathbf r}\,dt$;
quantum mechanically,
$P_{\rm bound}(t)=\sum_{nlm}|\langle nlm|\psi(t)\rangle|^{2}$ with
amplitudes $c_f(t)$ as above.

\subsection*{6. Worked example: capture by one drag pass through a gas envelope}

\paragraph{Model.}
A small body (radius $a$, density $\rho_b$, mass $m=\tfrac43\pi a^{3}\rho_b$) arrives at a planet
(mass $M$, radius $R$) with relative speed $v_\infty$ and impact parameter $b$.
It passes periapsis $r_p$ inside an extended gas envelope of outer radius $r_{\rm atm}$.
Assumptions: a two-body hyperbola before the encounter, a single \emph{impulsive}
drag pass at periapsis through a slab of constant column density $\Sigma=\int\rho_g\,ds$,
and a constant drag coefficient $C_D$.

\paragraph{Drag loss.} With $m\,v\,\dfrac{dv}{ds}=-\tfrac12 C_D A\rho_g v^{2}$ and $A=\pi a^{2}$,
\begin{equation}
v_{\rm out}=v_{\rm in}\,e^{-\kappa},\qquad
\kappa=\frac{C_D A\,\Sigma}{2m}=\frac{3\,C_D\,\Sigma}{8\,\rho_b\,a}.
\end{equation}
Smaller bodies have larger $\kappa$ ($\kappa\propto 1/a$), so they are captured more easily.

\paragraph{Capture condition.}
The periapsis speed is $v_p^{2}=v_\infty^{2}+2GM/r_p$. The specific energy after the pass is
\begin{equation}
\varepsilon'=\tfrac12 v_p^{2}\,e^{-2\kappa}-\frac{GM}{r_p}.
\end{equation}
Requiring $\varepsilon'<0$ gives
\begin{equation}
v_\infty^{2}<v_{\rm esc}^{2}(r_p)\,\big(e^{2\kappa}-1\big)
\quad\Longleftrightarrow\quad
r_p<r_c\equiv r_g\big(e^{2\kappa}-1\big),\qquad
r_g\equiv\frac{2GM}{v_\infty^{2}}.
\end{equation}

\paragraph{Cross-section with gravitational focusing.}
Conservation of $L=b\,v_\infty=r_p v_p$ gives
\begin{equation}
b^{2}(r_p)=r_p^{2}\Big(1+\frac{2GM}{r_p v_\infty^{2}}\Big)=r_p^{2}+r_g\,r_p .
\end{equation}
Capture occurs for $R\le r_p\le r_{\max}=\min(r_{\rm atm},r_c)$ (smaller $r_p$ means collision), so
\begin{equation}
\sigma_{\rm cap}=\pi\big[b^{2}(r_{\max})-b^{2}(R)\big]
=\pi\Big[(r_{\max}^{2}-R^{2})+r_g\,(r_{\max}-R)\Big],
\qquad
\Gamma=n_b\,\sigma_{\rm cap}\,v_\infty .
\end{equation}
Capture is impossible if $\kappa<\kappa_{\min}=\tfrac12\ln(1+R/r_g)$. The result is limited
by the envelope size if $\kappa\ge\tfrac12\ln(1+r_{\rm atm}/r_g)$.

\paragraph{Numbers (Jupiter).}
\begin{itemize}
  \item $GM=1.267\times10^{17}\,\mathrm{m^{3}\,s^{-2}}$, $R=7.149\times10^{7}\,\mathrm{m}$,
        $r_{\rm atm}=10R=7.149\times10^{8}\,\mathrm{m}$;
  \item $v_\infty=1\,\mathrm{km\,s^{-1}}$, which is comparable to $n r_H\approx0.9\,\mathrm{km\,s^{-1}}$
        with $r_H\approx5.3\times10^{10}\,\mathrm{m}$ and $n=1.68\times10^{-8}\,\mathrm{s^{-1}}$;
  \item $r_g=2GM/v_\infty^{2}=2.534\times10^{11}\,\mathrm{m}$, and $v_{\rm esc}(R)=59.5\,\mathrm{km\,s^{-1}}$;
  \item body: $a=1\,\mathrm{km}$, $\rho_b=2000\,\mathrm{kg\,m^{-3}}$, $C_D=1$,
        so $\kappa=1.875\times10^{-7}\,\Sigma$ with $\Sigma$ in $\mathrm{kg\,m^{-2}}$;
  \item thresholds: $\kappa_{\min}=1.41\times10^{-4}$ and $\kappa_{\rm env}=1.41\times10^{-3}$.
\end{itemize}

\begin{center}
\begin{tabular}{|c|c|c|c|c|}
\hline
$\kappa$ & $\Sigma\ (\mathrm{kg\,m^{-2}})$ & $r_{\max}$ & $\sigma_{\rm cap}\ (\mathrm{m^{2}})$ & $b_{\max}/r_H$\\
\hline
$10^{-1}$ & $5.3\times10^{5}$ & $r_{\rm atm}=7.15\times10^{8}\,\mathrm{m}$ (envelope-limited) & $5.1\times10^{20}$ & $0.25$\\
$10^{-3}$ & $5.3\times10^{3}$ & $r_c=5.07\times10^{8}\,\mathrm{m}\ (7.1\,R)$ (drag-limited) & $3.5\times10^{20}$ & $0.21$\\
$10^{-4}$ & $5.3\times10^{2}$ & $r_c<R$ & $0$ (collision only) & ---\\
\hline
\end{tabular}
\end{center}

The geometric area $\pi r_{\rm atm}^{2}=1.6\times10^{18}\,\mathrm{m^{2}}$ is enlarged by the focusing factor
$1+v_{\rm esc}^{2}(r_{\rm atm})/v_\infty^{2}\approx 355$, so the first row is
$\sigma_{\rm cap}\approx3\times10^{2}\,\pi r_{\rm atm}^{2}$ once the exclusion of $r_p<R$ is accounted for.

\paragraph{Resulting orbit ($\kappa=0.1$, $r_p=r_{\rm atm}$).}
\begin{align}
v_p^{2}&=v_\infty^{2}+\frac{2GM}{r_p}=3.554\times10^{8}\,\mathrm{m^{2}s^{-2}},\\
\varepsilon'&=\tfrac12 v_p^{2}e^{-0.2}-\frac{GM}{r_p}=-3.17\times10^{7}\,\mathrm{J\,kg^{-1}}<0,\\
a_{\rm sat}&=-\frac{GM}{2\varepsilon'}=2.0\times10^{9}\,\mathrm{m}\approx28\,R,\qquad
e=1-\frac{r_p}{a_{\rm sat}}\approx0.64,\\
r_a&=2a_{\rm sat}-r_p=3.3\times10^{9}\,\mathrm{m}\approx46\,R,\qquad
T=2\pi\sqrt{a_{\rm sat}^{3}/GM}\approx18\ \mathrm{d}.
\end{align}
Here $a_{\rm sat}/r_H\approx0.04$, so the captured orbit lies deep inside the Hill sphere.
The final state is bound ($E_{12}<0$) because the gas removed $\Delta E=-\tfrac12 m v_p^{2}(1-e^{-2\kappa})$.

\paragraph{Validity.}
The formulas are two-body results, accurate for $b_{\max}\ll r_H$. Here $b_{\max}/r_H\approx0.2$--$0.25$ is only
marginal, and the true cross-section should be obtained from Hill-problem integrations
with $C_J$ tracked. A real envelope has $\Sigma=\Sigma(r_p)$ rather than a uniform slab,
which replaces $r_{\max}=\min(r_{\rm atm},r_c)$ by the solution of $r_p=r_c(\kappa(r_p))$.
Later passes keep shrinking $a_{\rm sat}$, so this single-pass result is a first estimate.
\end{document}
